\documentclass[11pt]{article}

\usepackage[margin=1.05in]{geometry}
\usepackage[T1]{fontenc}
\usepackage{mathpazo}
\usepackage[scaled=0.92]{helvet}
\usepackage{courier}
\usepackage{amsmath,amssymb,mathtools,bm}
\usepackage{booktabs,array,multirow}
\usepackage{enumitem}
\usepackage{xcolor}
\usepackage{graphicx}
\usepackage{tikz}
\usepackage{tcolorbox}
\usepackage{titlesec}
\usepackage{caption}
\usepackage{microtype}
\usepackage{needspace}
\usepackage[hidelinks]{hyperref}

\usetikzlibrary{positioning,arrows.meta,calc}
\tcbuselibrary{skins,breakable}

% ---------------------------------------------------------------- palette
\definecolor{navy}{RGB}{25,55,92}
\definecolor{muted}{RGB}{90,96,105}
\definecolor{rule}{RGB}{200,201,196}
\definecolor{okgreen}{RGB}{0,110,60}
\definecolor{specblue}{RGB}{42,120,214}
\definecolor{newmag}{RGB}{160,45,120}
\definecolor{openamber}{RGB}{170,110,0}
\definecolor{limred}{RGB}{175,45,45}
\definecolor{tint}{RGB}{248,248,246}
\definecolor{lg}{RGB}{175,177,172}

% ---------------------------------------------------------------- headings
\titleformat{\section}{\normalfont\sffamily\large\bfseries\color{navy}}{\thesection}{0.7em}{}
\titleformat{\subsection}{\normalfont\sffamily\normalsize\bfseries\color{navy}}{\thesubsection}{0.6em}{}
\titleformat{\part}[display]
  {\normalfont\sffamily\Large\bfseries\color{navy}\filcenter}
  {\MakeUppercase{\partname~\thepart}}{0.4em}{}
\titlespacing*{\section}{0pt}{1.5em}{0.6em}
\titlespacing*{\subsection}{0pt}{1.2em}{0.45em}

% ---------------------------------------------------------------- macros
\newcommand{\R}{\mathbb{R}}
\newcommand{\Prob}{\mathbb{P}}
\newcommand{\Ind}{\mathbf{1}}
\newcommand{\dhat}{\hat{d}}
\newcommand{\pihat}{\hat{\pi}}
\newcommand{\nhat}{\hat{n}}
\newcommand{\chat}{\hat{c}}
\newcommand{\that}{\hat{t}}
\newcommand{\Dhat}{\hat{\Delta}}

\newcommand{\provtag}[2]{{\footnotesize\sffamily\textcolor{#1}{[#2]}}}
\newcommand{\VALID}{\provtag{okgreen}{VALIDATED}}
\newcommand{\SPEC}{\provtag{specblue}{SPECIFIED}}
\newcommand{\NEWT}{\provtag{newmag}{NEW}}
\newcommand{\OPENT}{\provtag{openamber}{OPEN}}
\newcommand{\LIMIT}{\provtag{limred}{LIMIT}}

\newcommand{\why}[1]{\par\smallskip\noindent\textcolor{muted}{\small\itshape #1}\par\smallskip}

\newtcolorbox{claimbox}[1][]{
  enhanced, breakable, colback=tint, colframe=navy, boxrule=0pt,
  leftrule=2.5pt, arc=0pt, outer arc=0pt,
  left=12pt, right=12pt, top=9pt, bottom=9pt, #1}

\newtcolorbox{resultbox}[1][]{
  enhanced, breakable, colback=okgreen!4, colframe=okgreen!55, boxrule=0pt,
  leftrule=2.5pt, arc=0pt, outer arc=0pt,
  left=12pt, right=12pt, top=9pt, bottom=9pt, #1}

\newtcolorbox{limitbox}[1][]{
  enhanced, breakable, colback=limred!4, colframe=limred!55, boxrule=0pt,
  leftrule=2.5pt, arc=0pt, outer arc=0pt,
  left=12pt, right=12pt, top=9pt, bottom=9pt, #1}

\newtcolorbox{codebox}[1][]{
  enhanced, breakable, colback=tint, colframe=rule, boxrule=0.5pt,
  arc=2pt, left=10pt, right=8pt, top=7pt, bottom=7pt,
  fontupper=\footnotesize\ttfamily, #1}

\newtcolorbox{notebox}[1][]{
  enhanced, breakable, colback=white, colframe=rule, boxrule=0.5pt,
  arc=2pt, left=10pt, right=10pt, top=7pt, bottom=7pt,
  fontupper=\small, #1}

\captionsetup{font=small,labelfont={sf,bf,color=navy}}
\setlist[itemize]{leftmargin=1.3em,itemsep=2pt,topsep=4pt}
\setlist[enumerate]{leftmargin=1.6em,itemsep=2pt,topsep=4pt}
\renewcommand{\arraystretch}{1.15}

\hypersetup{
  pdftitle={Experiment Designs and Pre-Registration: Experiments S and P},
  pdfauthor={Robert Romani},
  pdfsubject={Pre-registration, 13 August 2026},
}


\title{\textcolor{navy}{\sffamily Experiment Designs and Pre-Registration}\\[6pt]
\large\normalfont Experiment S (scaffold characterisation) and Experiment P (pooled recovery)}
\author{Robert Romani}
\date{August 13, 2026}

\begin{document}
\maketitle

\begin{abstract}
\noindent This document is a \textbf{pre-registration}. It states the design, the metrics, and
the predicted outcomes of Experiments S and P \emph{before either is run}, so that the results
can be reported as hits and misses rather than as a narrative assembled afterwards. The project
has twice benefited from doing this: two of the four future-work propositions were falsified and
the repairs became the backbone of the method, and the U-shape prediction for orientation error
was falsified on 13~August in a way that inverted a design axis.

\medskip
\noindent Both experiments previously existed only as prose in \S8 and \S9 of
the 31~July
detection-phase working document. This note extracts them, updates them
for what has been measured since, and records what would count as a failure.
\end{abstract}

\noindent\textbf{\sffamily\small\textcolor{navy}{Provenance tags.}}
\smallskip

{\small
\begin{tabular}{@{}ll@{}}
\VALID & measured against ground truth\\
\SPEC & designed, not yet run\\
\NEWT & changed or added since the 31 July design\\
\OPENT & acknowledged gap\\
\end{tabular}}

\vspace{0.4em}

\section{Why these two, and in this order}

The method note (\texttt{boundary\char`_recovery\char`_v5}) sets out a four-step critical path. Steps 1
and 2 are done:

\begin{itemize}
\item \textbf{Step 1} --- the distance estimator $\dhat = -\sigma\Phi^{-1}(\pihat)$ recovers true
distance at $r = 0.996$, median relative error $-0.92\%$, flat across a $10\times$ noise sweep.
\VALID
\item \textbf{Step 2} --- the separability-derived normal carries \emph{no detectable bias} in
any of 20 cells; pooled orientation error falls below $3^\circ$ at $N = 100$ down to
$\Delta/\tau = 1.5$. \VALID
\end{itemize}

\noindent Experiment P is step 3. Experiment S is not on that path at all --- it characterises
the \emph{real} target, and every power number in the project to date is against planted
synthetic gates that we designed ourselves.

\textbf{S runs first, and specifically S5 and S1.} They are cheap, and between them they decide
what the rest of the summer is for. If the scaffold's jump is $3\text{--}5\sigma$, detection
power is not a live problem and the binding constraints are coverage and the on-manifold
tension. And if the OOD detector does not recognise our probe geometry, the paper's claim
changes from ``we detect the switch'' to ``we are immune to it'' --- which reframes everything
upstream.

\part{Experiment S --- characterising the actual target}

\section{Premise and build}

\begin{claimbox}
\noindent Every power figure in this project is measured against gates we planted. Measure what
the real attack induces \emph{before} optimising further against synthetic ones.
\end{claimbox}

\noindent The target is the Slack et al.\ (AIES 2020) scaffolding attack: a biased model $f$
that uses a sensitive attribute, an innocuous $\psi$ that uses an uncorrelated synthetic
feature, and an OOD detector --- a random forest trained to separate real data from
LIME-perturbed samples --- routing between them.

\paragraph{Audit the score, not the label.} The framework presumes a continuous response. The
housing scaffold in \texttt{src/attacks/}\allowbreak\texttt{scaffold.py} returns continuous predictions, so this is
satisfied. If a dataset exposes only hard labels, stop and record that a discrete-response
Stage~3 variant is a prerequisite rather than a tweak.

\paragraph{Data availability, checked 13 August.} \NEWT\ The \textbf{housing scaffold is
runnable now}: \texttt{seattle\char`_sample\char`_3k.csv} and the demographics join are present, the
scaffold and both OOD detector constructors are in \texttt{src/attacks/}, and
\texttt{ood\char`_detector.}\allowbreak\texttt{predict\char`_proba} is directly callable. The \textbf{COMPAS, German Credit
and Communities-and-Crime datasets are absent} --- \texttt{papers/} contains only
\texttt{pexsbm}, not the Slack repository, and only downstream result directories survive. S
therefore runs on housing first; the other three require re-obtaining the Slack repository and
are deferred, not dropped.

\section{The five measurements}

\subsection*{S1 --- what is the scaffold's $\Delta/\sigma_{\text{resid}}$?}

Sample point pairs straddling the OOD decision boundary at matched covariate distance; measure
the jump in audited output relative to the LTS residual scale at the same anchors.

This is the single number that determines whether detection power is a live problem at all.
Reference points from the measured floor: the dip has \emph{zero} power below $\approx 2\sigma$
(the mixture is literally unimodal there), and the equal-variance LRT reaches $0.42$--$0.54$ at
$1.5\sigma$ and $1.00$ at $2.5\sigma$.

\subsection*{S5 --- what does the detector call \emph{our} probes?}

Generate the audit's probe cloud at real anchors and ask the OOD detector to classify it.

The audit's perturbation is not LIME's. LIME perturbs every feature, ambiently, with unit
Gaussian noise; the detector is trained on exactly that. Ours is isotropic within a
$\hat d$-dimensional tangent frame at small radius. \textbf{The detector was never trained
against our probe geometry.}

\subsection*{S2, S3, S4}

\begin{itemize}
\item \textbf{S2 --- boundary accessibility.} Fraction of real held-out points lying within
$r_{\text{hi}}$ of the OOD boundary. This is coverage for the real target, as opposed to
coverage in a simulation where anchors are placed by design.
\item \textbf{S3 --- density-filter tension.} For anchors whose ball crosses the OOD boundary,
record the retention $\rho$.
\item \textbf{S4 --- detector leakage.} Fraction of real points the OOD detector misclassifies.
This is the on-manifold region where crossing is possible without leaving support.
\end{itemize}

\begin{limitbox}
\noindent\textbf{S3 is an explicit falsification condition.} If $\rho$ falls below the abstention
threshold at essentially every anchor whose ball crosses the OOD boundary, then \texttt{A3} fires
wherever the signal is and \textbf{the method abstains on its own target}. Report that as a
falsification, not as a tuning problem.
\end{limitbox}

\section{The pre-registered trichotomy}

S5 determines which of three regimes the scaffold falls into. All three are reportable; two are
wins.

\begin{table}[h]
\small
\begin{tabular}{@{}cp{0.40\textwidth}p{0.42\textwidth}@{}}
\toprule
& \textbf{\sffamily What S5 shows} & \textbf{\sffamily What it means}\\
\midrule
(i) & probes are classified as real; we query $f$ throughout the ball
    & No gate to detect --- but our attributions are of the \emph{biased} model.
      \textbf{The scaffold is defeated by evasion rather than by detection.} A paper result,
      differently framed.\\
(ii) & probes straddle the OOD boundary
    & Gate detected. The method works as designed on the real target.\\
(iii) & probes are classified OOD; we see only $\psi$
    & Fooled. The density filter is supposed to prevent this --- verify that it does.\\
\bottomrule
\end{tabular}
\end{table}

\paragraph{Prediction, stated before running.} \NEWT\ \textbf{Regime (i).} The LIME detector is
trained on ambient $N(0, 0.3)$ perturbations across \emph{all} features; our probes live in a
$\hat d$-dimensional tangent frame with zero component off it. Different geometry entirely.
Concretely: \textbf{we predict the detector classifies $> 80\%$ of our probes as real}, and that
the audited response is $f$ rather than $\psi$ over the great majority of each probe ball.

\paragraph{The structural reading, if (i) holds.} Regimes (i) and S3 are two faces of one fact:
the gate fires off-manifold, so an on-manifold auditor either never sees it (immune) or must
leave the manifold to see it --- at which point the density filter rejects the probes and the
claim degrades from \emph{a gate fires on deployed inputs} to \emph{a gate exists somewhere in
input space}. Stating that cleanly is arguably a stronger contribution than a detection number:
\textbf{scaffolding attacks hide in the region an on-manifold audit is definitionally excluded
from.}

\part{Experiment P --- pooled recovery}

\section{Design principle}

\begin{claimbox}
\noindent False positives and true positives differ \emph{geometrically}, not merely in
strength. A real boundary produces a consistent normal across every anchor touching it; a
spurious detection produces a random one. Pooling on geometric agreement suppresses false
positives faster than it accumulates true ones --- which is what is wanted when per-anchor power
is around $0.5$.
\end{claimbox}

\noindent Two corollaries. \textbf{Do not threshold per anchor before pooling}: at power $0.5$
that discards half the true signal. And the figure of merit becomes the number of anchors
required to recover the boundary, not per-anchor power.

\section{The per-anchor primitive}

Computed for \textbf{every} anchor, including those that do not fire.

\begin{codebox}
\begin{verbatim}
1. 2-GMM EM on the LTS residuals            -> responsibilities gamma
2. linear discriminant on (z, gamma)        -> normal nu, offset b
   *** fit on ALL probes, not the LTS inlier set ***          [NEW]
3. lift to ambient:  n_a = U nu / ||U nu||
   oriented so the higher-mean component is on the + side
4. signed offset:    c_a = n_a . x_a + b / ||nu||
5. weight:           w_a = LRT statistic, NOT thresholded

   each anchor emits (n_a, c_a, d_hat_a, Delta_hat_a, w_a)
\end{verbatim}
\end{codebox}

\noindent Steps 1--3 are \VALID\ as of 13~August. Step 2's ``all probes'' qualifier is
\NEWT: median orientation error is $5.8^\circ$ fitting on every probe against $21.9^\circ$ on
the LTS inlier set, because the trim removes points that are predominantly the minority branch
--- exactly the branch carrying the crossing signal. Soft responsibilities also beat hard
$\gamma > 0.5$ labels ($5.8^\circ$ against $7.6^\circ$).

\section{Estimators and null}

\begin{table}[h]
\small\centering
\begin{tabular}{@{}cp{0.52\textwidth}p{0.28\textwidth}@{}}
\toprule
& \textbf{\sffamily Estimator} & \textbf{\sffamily Uses geometry}\\
\midrule
A & Fisher/Stouffer combination of per-anchor $p$-values & no --- baseline\\
B & threshold on $w_a$, then cluster $(n,c)$ in projective space (antipodal identification),
    requiring $|\cos(n_a,n_b)| > 1-\epsilon$ and $|c_a - c_b| < \delta$ & partly\\
C & $w$-weighted mode-seeking (mean-shift or Hough voting) over \emph{all} anchors, no
    thresholding & fully\\
\bottomrule
\end{tabular}
\end{table}

\paragraph{Null.} Permute the residual-to-point assignment within each anchor and rerun the
pipeline from step 1. This randomises the normals while preserving the marginal residual
distributions. $\ge 500$ permutations.

\section{Generative setting} \NEWT

Flat intrinsic $d = 2$ manifold in ambient $D = 20$; a single planted hyperplane; the tangent
frame supplied, per the Stage-0 deferral.

\begin{table}[h]
\small\centering
\begin{tabular}{@{}lll@{}}
\toprule
\textbf{\sffamily Axis} & \textbf{\sffamily 31 July design} & \textbf{\sffamily Revised, 13 August}\\
\midrule
$\Delta/\sigma_{\text{resid}}$ & $\{1.0, 1.5, 2.5, 5.0\}$ & unchanged\\
$\pi$ & $\{0.10, 0.25, 0.50\}$ & $\bm{\{0.05,\ 0.10,\ 0.20\}}$\\
$N$ anchors & $\{5,10,25,50,100,200\}$ & unchanged\\
anchor distance $/\,r_{\text{hi}}$ & $\{0.2, 0.5, 0.9, 1.5\}$
  & \textbf{shell} $\bm{d/\sigma \in [1.28, 1.64]}$, plus a non-crossing control\\
\bottomrule
\end{tabular}
\end{table}

\paragraph{Why the two axes moved.} Orientation error is \textbf{monotone increasing in $\pi$},
not U-shaped as originally predicted: per-anchor standard deviation runs $3.9^\circ$ at
$\pi = 0.05$ to $47.5^\circ$ at $\pi = 0.50$ (at $\Delta/\tau = 5$). Fewer crossers, but they
sit further out along the normal, and the separation of class-conditional means along $n$
\emph{rises} as $\pi$ falls --- $2.17\sigma$ at $\pi = 0.05$ against $1.60\sigma$ at
$\pi = 0.50$. At $\pi \to 0.5$ there is additionally no majority branch for LTS to commit to.

Since $\pi \le 0.10$ means $d \ge 1.28\sigma$ and the minimum-mass rule caps $\pi \ge 0.05$
at $d \le 1.64\sigma$, the orientation-useful anchors occupy a \textbf{shell}. The original
distance axis sampled the wrong side of it.

\begin{limitbox}
\noindent\textbf{A deployment consequence P must measure.} That shell is roughly 22\% of the
anchors that can detect anything at all, which are themselves 15--20\% of randomly placed
anchors --- so on the order of \textbf{3--4\% of anchors in deployment}. Pooling is weighted
rather than thresholded, so higher-$\pi$ anchors still contribute, but $N_{95}$ under
\emph{random} placement may be far larger than $N_{95}$ under placement by design.
\textbf{Report both.}
\end{limitbox}

\section{Metrics}

\begin{enumerate}
\item $N_{50}$ and $N_{95}$ per estimator per $\Delta$, \textbf{under both by-design and random
anchor placement}. \NEWT\ Headline.
\item Estimator ranking.
\item $N \mapsto$ orientation error and $N \mapsto |\chat - c^\star|$. The first is partly
previewed (\S1); the second is untested.
\item False-positive rate under the permutation null.
\end{enumerate}

\noindent Secondary: two parallel boundaries at separation $\in \{0.5, 2, 5\}\cdot r_{\text{hi}}$
--- does mode-seeking resolve two modes or merge them? This sets the resolution limit for telling
one gate from two. And a curved boundary (the deep-gate preimage): pool locally, then link, and
report where the local-hyperplane approximation fails as a function of curvature $\times\,
r_{\text{hi}}$.

\section{The no-information guard} \NEWT

\begin{limitbox}
\noindent Outside the operating region --- $\Delta/\tau = 1.0$ at any $\pi$, or
$\pi \gtrsim 0.35$ --- the per-anchor rotation has standard deviation $\approx 52^\circ$, which
is exactly what uniformly random directions give. Pooling uninformative \emph{axial} estimates
does not converge to the truth and does not converge to nothing: the mean direction of $N$
random axes is itself uniformly random, so the pooled estimate lands on an arbitrary direction
with an interval that \emph{tightens} around it.
\end{limitbox}

\noindent \textbf{Pooling can therefore be confidently wrong in a way that per-anchor abstention
never is.} This is the second failure mode of that character in the method --- the first being
the $L < r_{\text{hi}}$ degeneracy, where a ball spanning several routing cells reads as
heteroscedastic noise and Stage~0 is blind to it.

P must include an explicit check that the anchor set carries information before it reports. The
permutation null is the natural instrument, since a genuinely uninformative anchor set should be
statistically indistinguishable from a permuted one. \textbf{Designing and calibrating that
guard is part of P, not a follow-up.}

\part{Pre-registered predictions}

\section{Stated before running; report hits and misses}

\begin{table}[h]
\small
\begin{tabular}{@{}cp{0.36\textwidth}p{0.44\textwidth}@{}}
\toprule
\textbf{\sffamily \#} & \textbf{\sffamily Prediction} & \textbf{\sffamily What falsifies it}\\
\midrule
S-1 & The OOD detector classifies $> 80\%$ of our tangent-frame probes as \emph{real} ---
      trichotomy regime (i). & $< 80\%$; or probes are called OOD, which is regime (iii).\\
S-2 & The scaffold's $\Delta/\sigma_{\text{resid}}$ exceeds $2.5$, so detection power is not the
      binding constraint. & A measured ratio below $\approx 1.5$, which puts the real target
      under the floor.\\
S-3 & Retention $\rho$ at boundary-crossing anchors falls below $\rho_{\min}$, so the audit
      abstains where the signal is. & $\rho$ comfortably above threshold --- in which case the
      on-manifold tension is not binding here.\\
P-1 & Estimator ranking $C > B > A$ at every $\Delta$, with the $C - B$ gap largest at
      $\Delta = 1.5\sigma$. & Any reordering; in particular $B \ge C$ would mean thresholding
      costs nothing.\\
P-2 & A real gate produces a \textbf{single mode} in the pooled offset $c$; resonant curvature
      produces a \textbf{comb} of modes spaced by its wavelength $\ell$. & A comb for a gate, or
      a single mode for resonant curvature --- either would put smooth rejection back on the
      per-anchor decision rules.\\
P-3 & $N_{95}$ under random anchor placement exceeds $N_{95}$ under by-design placement by
      roughly the inverse shell fraction ($\sim 20\text{--}30\times$). & A much smaller
      penalty, meaning higher-$\pi$ anchors contribute more than the per-anchor numbers imply.\\
P-4 & Orientation error scales like $\sqrt{d/m}$ in \emph{intrinsic} dimension, so it degrades
      at $d = 4, 8$. \OPENT\ (not run at $d = 2$) & Error flat in $d$, which would extend E1's
      dimension-independence result from ambient to intrinsic dimension.\\
\bottomrule
\end{tabular}
\end{table}

\section{Order of execution}

\begin{enumerate}
\item \textbf{S5, then S1}, on the housing scaffold. Cheap; determines the framing.
\item \textbf{S3}, since it is the explicit falsification condition and follows directly from
S5's machinery.
\item \textbf{Experiment P}, with the revised axes and the no-information guard.
\item S2 and S4; the COMPAS / German / Communities replication once the Slack repository is
re-obtained.
\end{enumerate}

\noindent Results are reported against the table in \S13 --- hits and misses, in that form.

\end{document}
